% =============================================================================
% Chapter 5 --- Empirical evidence (Ju & Aral 2026)
% =============================================================================
\chapter{Empirical evidence: Ju \& Aral (2026)}
\label{ch:empirical}

\section{Source document}

\begin{description}
  \item[Title.] Collaborating with AI Agents: A Field Experiment on
    Teamwork, Productivity, and Performance.
  \item[Authors.] Harang Ju (Johns Hopkins Carey) and Sinan Aral (MIT
    Sloan).
  \item[Venue.] arXiv:2503.18238v3 [cs.CY], 5 February 2026.
  \item[Pages.] 59.
  \item[Source.] \corpus{source/2503.18238v3.pdf}
    (sha256 \texttt{b680f584\ldots cd7877}).
\end{description}

\section{Study design}

The authors built the \emph{Pairit} platform and randomly assigned 2{,}234
participants to human-human and human-AI teams producing 11{,}024 ads
for a think tank. Outputs were rated by humans and field-tested on X,
which garnered roughly five million impressions. The design is --- for
this corpus --- an unusually strong source of evidence about how
human-agent collaboration actually affects output quality, homogeneity,
and team dynamics.

\section{Findings relevant to governance}

\begin{enumerate}
  \item \textbf{Jagged productivity frontier.} Human-AI teams produced
        50\% more ads per worker and higher \emph{text} quality;
        human-human teams produced higher \emph{image} quality. Different
        sub-tasks need different evaluation metrics.
  \item \textbf{Delegation shifts.} Participants delegated 17\% more work
        to AI agents than to human partners and performed 62\% fewer
        direct text edits when working with AI.
  \item \textbf{Diversity collapse.} Human-AI outputs were more
        self-similar. For creative or generative use cases this is a
        first-order governance risk.
  \item \textbf{Communication restructure.} Human-AI collaboration was
        significantly more task-oriented (+25\% task messages, -18\%
        interpersonal).
  \item \textbf{AI-awareness moderation.} Participants who correctly
        identified they were paired with AI delegated more and were more
        task-oriented --- the effect is moderated by knowing the
        collaborator is an agent.
\end{enumerate}

\section{Requirements extracted}

\begin{table}[h]
\centering
\small
\begin{tabularx}{\textwidth}{p{3.2cm}p{1.4cm}p{1.6cm}X}
\toprule
\textbf{Req.\ id} & \textbf{Section} & \textbf{Obligation} & \textbf{Normative text (paraphrased)} \\
\midrule
REG-JU-ARAL-3-001 & 3 & descriptive & Evaluate per-task; do not generalise a single productivity metric. \\
REG-JU-ARAL-4-001 & 4 & should & High delegation + low direct-edit rate is a proxy for automation-bias risk; measure it. \\
REG-JU-ARAL-5-001 & 5 & should & Add an output-diversity metric for creative-output use cases. \\
\bottomrule
\end{tabularx}
\caption{Observational claims treated as requirements.}
\label{tab:jural-reqs}
\end{table}

\section{Implementation in the corpus}

\begin{itemize}
  \item REG-JU-ARAL-3-001 is implemented by this chapter; it reminds
        reviewers to look at the per-task KPIs defined in
        \corpus{../tb/structured/kpis/tb-kpis.json} rather than a single
        blended productivity number.
  \item REG-JU-ARAL-4-001 maps to the override-rate and review-duration
        KPIs in the TB corpus. When override rates fall well below
        historical baseline and review duration drops in lockstep, that
        is the signal to re-audit.
  \item REG-JU-ARAL-5-001 is marked \emph{not\_applicable} for the
        trial-balance workflow (deterministic review, single narrative
        per legal entity). It will become applicable when the corpus
        expands to generative use cases; we record it rather than omit it
        so that the framework surfaces the need at that time.
\end{itemize}
\clearpage
